\documentclass[12pt,a4paper]{article}

% ── Packages ──────────────────────────────────────────────────
\usepackage[T1]{fontenc}          % Correct font encoding for output
\usepackage[utf8]{inputenc}       % Allows UTF-8 characters (optional in modern LaTeX)
\usepackage{microtype}            % Improves typography subtly
% \usepackage{palatino}             % Palatino text font
% \usepackage{mathpazo}             % Palatino-compatible math font
\usepackage[margin=2.5cm]{geometry} % Page layout
\usepackage{setspace}             % Line spacing control
\usepackage{parskip}              % Adds space between paragraphs
\usepackage{xcolor}               % Colors for text and backgrounds
\usepackage[hidelinks]{hyperref}  % Clickable hyperlinks
\setlength{\parskip}{8pt plus 2pt minus 1pt}

% ── Typography ────────────────────────────────────────────────
\onehalfspacing


% Hyperref metadata
\hypersetup{
  pdftitle={Regulatory Texts and Case Law as Ground Truth in Emerging Domains},
  pdfauthor={Lalitha A R},
  pdfsubject={Research Methodology; Legal Informatics; Food Systems},
  pdfkeywords={regulatory delta; legal ground truth; food informatics; methods}
}


\title{Regulatory Texts and Case Law as Ground Truth\\ in Emerging Domains}

\author{Lalitha A R\\
\texttt{lalithaar.research@gmail.com}\\
Lead Researcher\\
Interdisciplinary Systems Research Lab (\textit{i}SRL)}

% ── Document ──────────────────────────────────────────────────
\begin{document}
\maketitle


\begin{abstract}
\noindent
\normalsize
In domains where no unified academic theory has yet consolidated, practitioners
face a calibration problem: against what baseline does one evaluate a classification
framework, a taxonomy, or a data model? This note argues that enacted law and
judicial decisions constitute a practically available, historically grounded source
of ground truth for such domains. We describe an approach---regulatory delta
analysis---that examines how legislation and case law shift across time, with the
aim of surfacing friction points, coordination patterns, and the constraints under
which each version of a rule was written. The approach is not a critique of
legislative bodies or courts; it is a method for reading the record they have
already produced. Limitations are considered throughout. The note is grounded in
applied work from the Indian Food Informatics Data (IFID) project but the
reasoning is intended to transfer to other domains with similar structural
properties.
\end{abstract}


\newpage
% ── 1. Definitions ──────────────────────────────────────────
\section{Definitions}

The terms below are used precisely throughout this note.

\textbf{Regulatory text.}\enskip A statute, rule, or formal regulation enacted by a
legislative or administrative body with binding legal authority. In the Indian
context, examples include the Food Safety and Standards Act, 2006, and the FSSAI
Labelling and Display Regulations, 2020.

\medskip
\textbf{Case law.}\enskip The body of judicial decisions---judgments handed down by
courts and tribunals---that interpret, apply, and sometimes settle conflicts
between regulatory texts. Case law acquires authority through the doctrine of
precedent and, in many common law systems including India's, through formal
hierarchical citation obligations.

\medskip
\textbf{Regulatory delta.}\enskip The set of meaningful changes between two enacted
versions of a regulatory text: additions, deletions, scope expansions,
definitional tightenings, and any shifts in the underlying policy stance. A delta
is not a simple diff; it is an interpretation of what changed and, where the
record permits, why.

\medskip
\textbf{Friction point.}\enskip A location in the regulatory record---a contested
statutory term, a split between regulators, a gap exploited by litigation---where
the interests or interpretations of two or more actors have visibly come into
tension. Friction points are distinct from errors; they are structurally revealing.

\medskip
\textbf{Ground truth (working definition).}\enskip For the purposes of this note, a
ground truth is a reference against which the outputs of an analytical model can be
tested. It does not imply perfect correctness; it implies that the reference has
been produced by a process that is independent of and prior to the analysis being
evaluated.

% ── 2. The calibration problem in emerging domains ──────────
\section{The Calibration Problem in Emerging Domains}

Research in well-consolidated fields benefits from a body of replication,
meta-analysis, and theoretical synthesis that can serve as a reference. A new
model or framework can be tested against this existing record. The food systems
informatics domain does not yet have this infrastructure in India. The field is
recent, the data are fragmented, and the constructs---what counts as the same
ingredient, how processing changes identity, how regional naming should relate to
regulatory naming---remain unsettled.

A comparable situation arises in any domain that is technically complex, involves
multiple institutional actors with partly overlapping mandates, and has developed
faster than the academic literature has been able to consolidate it. Infrastructure
regulation, environmental classification, digital governance, and traditional
medicine systems all share these properties to varying degrees.

In such conditions, one cannot simply look to an established canon of empirical
findings to validate a new framework. The question then is: what stable, legible,
independently produced record exists against which one can calibrate?

We propose that the regulatory and legal record fills this role, with specific
properties and specific limitations that are developed in the sections that follow.

% ── 3. Why law ──────────────────────────────────────────────
\section{Why Law and Case Law Can Function as Ground Truth}

Regulatory texts and judicial decisions share several properties that make them
useful as a reference in calibration contexts.

\textbf{They are the result of adversarial processes.} Legislation is typically
produced after consultation, lobbying, expert review, and political negotiation.
Judicial decisions are produced after argument by opposed parties, subject to
appeal, and written to justify a conclusion against the best available contrary
reading. Both processes are imperfect, but both are designed to surface objections.
The record they produce has, in a meaningful sense, survived challenge.

\textbf{They name conflicts directly.} Academic literature tends to report
consensus or to frame disagreement theoretically. Case law reports disagreement
factually: here are two parties, here is what they disputed, here is which
interpretation prevailed and why. The friction is the content of the document
rather than a subtext to be inferred.

\textbf{They are time-stamped.} Each regulatory text and each judgment carries a
date. This makes it possible to sequence the record chronologically and ask which
understanding of a concept was operative when.

\textbf{They are publicly accessible.} In India, statutory instruments are notified
in the Gazette of India. Judgments of the Supreme Court and High Courts are
published in official reporters and on government databases. The record is, in
principle, reachable by any researcher.

\textbf{They have institutional authority within their domain.} A food safety
regulation issued under the FSSAI Act is, for practical purposes, the definition
of the relevant concept for the actors it governs---manufacturers, importers,
inspection officers---regardless of whether a food scientist would agree with it.
When building systems that must operate in that legal environment, the legal
definition is not one input among many; it is a constraint.

None of these properties make the legal record infallible. Section~6 addresses
limitations. But they are sufficient to make law a productive starting point when
other reference points are unavailable.

% ── 4. Regulatory delta analysis ────────────────────────────
\section{Regulatory Delta Analysis as Method}

A regulatory delta is not simply a list of changes between two versions of a
statute or rule. It is an interpretation of those changes in light of the context
that produced them.

\subsection{Reading changes in context}

Laws are written under constraints. The constraints include the state of the
relevant industry at the time of drafting, the administrative capacity available to
enforce the rule, the incidents that made a regulatory response necessary, and the
political feasibility of various options. A later version of a rule almost always
looks more precise, more comprehensive, or more technically sophisticated than an
earlier one---not because earlier drafters were careless, but because they were
working with less data, less precedent, and a less developed field.

The appropriate stance when reading a delta is that of a scribe rather than a
judge: the task is to document what changed, to note what the earlier version could
not have anticipated, and to ask what new information or new pressure made the
change necessary. This is a different question from asking whether the earlier rule
was wrong.

Applied to the Indian food labelling context, the transition from the Food Safety
and Standards (Packaging and Labelling) Regulations, 2011 to the Food Safety and
Standards (Labelling and Display) Regulations, 2020 illustrates this clearly
\cite{ifid:regs2026}. The 2011 regulations were drafted at a point when
industrial food processing in India was still expanding rapidly and digital
traceability tools were not yet available to enforcement bodies. The 2020
regulations tightened allergen declarations, introduced structured front-of-pack
warnings, and prescribed naming conventions with greater specificity. Each of these
additions corresponds to a domain that developed, was observed, and was then
addressed. The delta reveals an institution processing experience and updating its
instrument accordingly.

\subsection{Identifying friction points through case law}

Where the regulatory text leaves ambiguity, the courts resolve it---and in doing
so, produce a record of where the ambiguity was, who held which interpretation, and
which reading eventually prevailed. This makes case law particularly useful for
identifying friction points that would not be visible from the statutory text alone.

The Supreme Court of India's January 2026 judgment in \textit{Commissioner of
Customs (Import) v.\ M/s Welkin Foods} illustrates this \cite{ifid:sc2026}. The
case concerned whether imported aluminium shelving should be classified as an
agricultural machine part or as an aluminium structure. The legal question was
narrow, but the Court's reasoning established a hierarchy for resolving
classification disputes in which statutory technical definitions take precedence
over common commercial understanding. This hierarchy had been contested in earlier
decisions and was now settled. For a food informatics system that must align with
Indian classification practice, this judgment is a material constraint---one that
would not have been legible from reading only the statutory text.

\subsection{Mapping coordination across institutional actors}

A regulatory domain typically involves multiple bodies with overlapping but
non-identical mandates. In Indian food systems, the relevant actors include the
Food Safety and Standards Authority of India (FSSAI), the Directorate General of
Foreign Trade (DGFT), the Central Board of Indirect Taxes and Customs (CBIC), and
the courts. These actors do not always interpret the domain identically, and their
instruments do not always align.

The legal record makes these relationships visible. Where one body's definition
conflicts with another's, there will typically be a judgment or a regulatory
amendment that resolves the conflict, defers it, or acknowledges it. Mapping these
interactions across time reveals not just what the current rule is, but how the
current rule came to be and which pressures it is still absorbing.

% ── 5. What this approach surfaces ──────────────────────────
\section{What This Approach Surfaces}

Regulatory delta analysis, applied systematically, tends to surface four types of
information that are difficult to obtain from other sources.

\textbf{Constraint archaeology.} Earlier versions of rules encode the constraints
that were operative when they were written. Identifying these constraints---and
asking whether they are still valid---can reveal where a regulatory framework is
load-bearing on an assumption that may no longer hold.

\textbf{Coordination mechanisms.} When two bodies with overlapping mandates produce
consistent rules over time, it is worth asking how that consistency is achieved.
The legal record will often contain evidence of formal coordination mechanisms,
mutual referencing, or the adoption of shared definitions.

\textbf{Friction without resolution.} Not all conflicts in the legal record are
resolved. Some cases are settled before judgment. Some regulatory ambiguities are
explicitly deferred. These unresolved tensions are as informative as the settled
ones: they mark the places where the system has not yet stabilised.

\textbf{Bias documentation.} Legal instruments are written by people operating in
institutional contexts, and they reflect the concerns, categories, and blind spots
of those contexts. A regulatory text that focuses on industrial food and does not
address traditional preparations is not neutral; it reflects what was legible and
politically salient at the time of drafting. Noting these asymmetries is part of
using the legal record honestly.

% ── 6. Limitations ──────────────────────────────────────────
\section{Limitations}

The approach described here has several limitations that must be held in mind.

\textbf{Law is not science.} A court may settle a dispute by choosing an
interpretation that is administratively convenient or politically feasible rather
than technically accurate. A regulatory definition may persist after the scientific
understanding of the relevant phenomenon has moved on. The legal record reflects
what was decided, not necessarily what was correct.

\textbf{Unenforced rules are not reliable evidence.} A statute that exists on paper
but is not enforced tells us something about legislative intent but little about
actual practice. The gap between enacted law and operational practice can be
substantial.

\textbf{The record is not complete.} Not all decisions are published. Not all
conflicts result in litigation. Regulatory negotiations that produce an amended
rule may leave no public trace of the original disagreement. The legal record
samples the domain rather than covers it.

\textbf{Jurisdiction specificity.} The regulatory architecture of one country or
regulatory system is not directly transferable to another. Insights from Indian
food law are not automatically generalisable to food systems in other jurisdictions,
though the structural properties of the method may transfer.

\textbf{Temporal lag.} Legislation and litigation are slow. The legal record may be
significantly behind the current state of the domain, particularly in fast-moving
technical fields. Using the legal record as ground truth requires acknowledging that
it may be calibrated to a version of the domain that no longer obtains.

These limitations do not disqualify the approach. They specify the conditions under
which it is and is not useful, and they indicate the supplementary sources---field
research, domain expert consultation, technical audits---that should accompany it.

% ── 7. Relationship to other methods ────────────────────────
\section{Relationship to Other Sources and Methods}

This note does not argue that the legal record should replace other methods of
establishing ground truth. It argues that the legal record is an underused source
that has specific properties making it productive in specific conditions: emerging
domains, multi-actor regulatory environments, and contexts where the gap between
legal definition and operational practice is itself an object of study.

The approach is most useful in combination with domain expert consultation, which
can identify where the legal record is silent or misleading; with empirical data
collection, which can reveal practice that diverges from legal prescription; and
with traditional academic literature, which provides theoretical frameworks for
interpreting what the legal record contains.

Academic literature is not deprecated here. The claim is narrower: in a domain
where the academic literature is still being assembled, the legal record is
available now, has been produced by adversarial processes, and carries authority
for the actors whose behaviour one is trying to understand or model. It is a
reasonable place to start.

% ── 8. Closing remarks ──────────────────────────────────────
\section{Closing Remarks}

Regulatory systems change. The record of that change is publicly available, time-stamped,
and produced by institutions that have observed the domain, absorbed feedback
from it, and updated their instruments accordingly. Reading that record carefully
is not an alternative to original research---it is a form of original research,
and one that takes the accumulated work of regulatory bodies and courts seriously
as evidence rather than setting it aside in favour of sources that may be more
recent but less tested.

The goal is not to judge who was right and who was wrong in any given dispute, or
whether a regulatory body made the best possible decision with the information
available. The goal is to understand where the system has been, where it has
strained, and what that history reveals about where it currently stands. That is a
question that the legal record is unusually well placed to answer.

\newpage
% ── Acknowledgments ─────────────────────────────────────────
\section*{Acknowledgments}

My deepest gratitude to Mr. Krishna, whose constancy forms the foundation upon which all my work, including this, quietly rests.

Salutations to the Goddess who dwells in all beings in the form of intelligence. I bow to her again and again.


This note draws on methods developed in the course of the Indian Food Informatics
Data (IFID) project at iSRL. The authors thank the researchers whose applied work
surfaced the need for explicit methodological documentation.

% ── References ──────────────────────────────────────────────
\begin{thebibliography}{9}
\setlength{\itemsep}{0.4em}

\bibitem{ifid:sc2026}
Lalitha, A.~R. (2026).\newblock {Indian Supreme Court Defines Hierarchical Classification for Food Products: Overruling Common Parlance Precedents}.
\newblock Interdisciplinary Systems Research Lab, February 2026.
\newblock \url{https://doi.org/10.5281/zenodo.18651645}.

\bibitem{ifid:regs2026}
Vukka, S.~N. \& Lalitha, A.~R. (2026).
\textit{Regulatory Delta of Food Labelling Laws in India: A Comparative
Analysis of the FSSAI 2011 and 2020 Regulations}.
Indian Food Informatics Data (IFID) Project,
Interdisciplinary Systems Research Lab.
\url{https://doi.org/10.5281/zenodo.18719394}

\bibitem{india:fssaact}
Government of India (2006).
\textit{Food Safety and Standards Act, 2006} (No.\ 34 of 2006).
Ministry of Health and Family Welfare.

\bibitem{india:welkin}
Supreme Court of India (2026).
\textit{Commissioner of Customs (Import) v.\ M/s Welkin Foods},
2026 SCC OnLine SC 27; 2026 INSC 19. January 6, 2026.

\end{thebibliography}
\end{document}
